% Source: 03 - Wed Sep 30 2026.pdf, page 8. Content remains in source order.
\sourcepage{03}{8}

\subsection{Reductions and Order}
\textbf{Theorem:}
Given $L_1,L_2\subseteq\bits^*$:
\[
(L_1\polyred L_2)\land(L_2\in\classP)\ \Rightarrow\ L_1\in\classP.
\]

If $L_1$ is reducible to $L_2$ and $L_2$ can be decided in poly-time, then $L_1$ can also be decided in poly-time

\textbf{Proof:} We have two hypotheses:

HP1: $L_1$ is poly-time reducible to $L_2$:
$L_1\polyred L_2:\ \exists f$ polynomially time computable (ptc) : $x\in L_1\Longleftrightarrow f(x)\in L_2$.

So we hypothesize that an algorithm $A_f$ exists that computes $f$ in polynomial time. Let $T_{A_f}(|x|)=O(|x|^{k_1})$, $k_1\geq0$.

HP2:
$L_2\in\classP:\ \exists A_{L_2}$ which decides $L_2$ in poly time. Let

$L_2$ is in P, so there exists an algorithm $A_{L_2}$ which decides $L_2$ in polynomial time. Let$T_{A_{L_2}}(|x|)=O(|x|^{k_2})$.

We can unite these two hypothesis into the following algorithm:

\textbf{$A_{L_1}(x)$}
\begin{tabbing}\quad\=\kill\> $y\leftarrow A_f(x)$\\\> return $A_{L_2}(y)$\end{tabbing}

So to solve instances of $L_1$, we first compute $y=f(x)$ and then we decide if $y\in L_2$. Since both f and $A_{L_2}$ are polynomial time computable, we can compute the running time of $A_{L_1}$:

Running time:
\[
T_{A_{L_1}}(|x|)=O(|x|^{k_1}+|y|^{k_2}).
\]
How large can $y$ be? $A_f$ may use all of its steps to create bits of $y$, thus
\[
|y|=O(|x|^{k_1}).
\]
Hence
\[
T_{A_{L_1}}(|x|)=O(|x|^{k_1}+|x|^{k_1k_2})=O(|x|^{k_3}),
\quad \text{where } k_3=\max\{k_1, k_1k_2\}.
\]
Which is still polynomial in $|x|$. 